\section{多源热点判别与正式结果审计}

\subsection{多源共同指向的操作定义}

赛题任务三要求从机会记录中识别由多个来源网站共同指向的热点方向。来源覆盖不能由记录条数替代，因此本文把主题覆盖的不同网站数记为 $s_c$，并以 $s_c\geq2$ 作为多源热点的操作判据。该判据只描述来源覆盖，不假定不同网站之间相互独立，也不排除同一事件在不同平台上的转载。对每个主题同时保存成员编号、来源名称和正文快照路径，便于核验是否存在重复传播。

正式结果包含 14 个展示项，其中 HDBSCAN 形成 2 个非噪声簇，12 条噪声记录按单记录展示项保留。C2 包含 22 条记录，覆盖 GitHub、V2EX 和量子位，满足多源热点定义。C1 包含 66 条记录，但全部来自 Hacker News，虽然规模较大，却不满足多源共同指向条件。其余 12 个展示项均为单记录、单来源，不能据此提出跨来源热点结论。

\begin{table}[H]
  \centering
  \small
  \caption{多源热点判别与聚类结构}
  \label{tab:multisource-audit}
  \begin{tabular}{lrrp{0.40\textwidth}l}
    \toprule
    主题 & 记录数 & 来源数 & 来源覆盖 & 多源热点 \ \\
    \midrule
    C2 & 22 & 3 & GitHub、V2EX、量子位 & 是 \\
    C1 & 66 & 1 & Hacker News & 否，单源对照 \\
    C3--C14 & 1 & 1 & 各自单一来源 & 否，单记录展示项 \\
    \bottomrule
  \end{tabular}
\end{table}

因此，C1 的价值分只表示单一平台内部的大规模记录集合在当前评分先验下的排序位置，不能写成多源热点。C2 的高分同时受到来源覆盖、记录规模、互动信号和时间新近性的影响，但其成员场景仍然异质，策划案只能把它组织为待验证的应用工作台方向。使用主题结构和来源覆盖共同观察热点，也与文献知识图谱研究将主题关系与关注方向结合分析的思路一致\cite{wang2023soilmining}。

\subsection{正式价值分与图表一致性}

正式结果使用四维评分：多源指向度、记录规模、互动信号和时间新近性，权重分别为 $0.40$、$0.30$、$0.15$ 和 $0.15$。正式 $result4.csv$ 的排序为 C2（81.52）、C1（53.31）、C14（39.12）、C13（39.10）、C12（39.08）、C11（33.93）、C10（32.37）、C8（26.71）、C9（26.20）、C7（24.90）、C6（23.90）、C4（23.49）、C3（23.44）和 C5（23.32）。图~\ref{fig:topicscore} 与表~\ref{tab:topic-scores} 均由同一正式主题运行目录生成，图中直接标出四维正式分值，避免使用早期三维评分结果。

\begin{table}[H]
  \centering
  \small
  \caption{正式主题价值分复核清单}
  \label{tab:formal-score-audit}
  \begin{tabular}{lrrrr}
    \toprule
    主题 & 记录数 & 来源数 & 多源热点 & 机会价值分 \\
    \midrule
    C2 & 22 & 3 & 是 & 81.52 \\
    C1 & 66 & 1 & 否 & 53.31 \\
    C14 & 1 & 1 & 否 & 39.12 \\
    C13 & 1 & 1 & 否 & 39.10 \\
    C12 & 1 & 1 & 否 & 39.08 \\
    C11 & 1 & 1 & 否 & 33.93 \\
    C10 & 1 & 1 & 否 & 32.37 \\
    其余 7 个单记录项 & 1 & 1 & 否 & 23.32--26.71 \\
    \bottomrule
  \end{tabular}
\end{table}

\section{置信度、阈值与结果表语义}

赛题表 2 的“机会置信度”定义为记录构成 AI 应用机会的可信度，数值越大，表示越支持机会判断。模型接口内部保留原始自评值 $q_i$，该值表示模型对当前四项判据的确定程度，尚未进行概率校准。为了使官方 CSV 的列语义与赛题一致，结果表输出机会方向分数

\begin{equation}
 p_i=\begin{cases}
 q_i,& y_i=1,\\
 1-q_i,& y_i=0.
 \end{cases}
\end{equation}

机会候选的程序阈值为 $t=0.5$。四项判据满足且 $q_i\geq t$ 时，记录进入原模型机会候选集合；命中排除项、缺少需求条件或低于阈值时，记录保持为非机会。$t=0.5$ 取未校准 0 到 1 自评区间的中点，作用是过滤低确定度候选，不被解释为经验校准概率。全量原模型输出得到 269 条候选机会，最终人工确认字段中保留 100 条机会。

在正式 $result2.csv$ 中，机会记录的 $p_i$ 范围为 0.78--0.99，中位数为 0.94；非机会记录的 $p_i$ 范围为 0--0.80，中位数为 0.02。少数非机会记录仍可能具有较高的机会方向分数，但由于四项条件未同时满足，硬判据优先于自评排序。这一设计同时保留了模型不确定性信息和赛题要求的二分类语义。

\section{补充文献与方法选择}

新增文献为方法选择提供了三个补充参照。数字心理健康平台研究使用 LDA、情绪和语义网络分析构造需求强度，说明主题频数本身不足以代表需求优先级\cite{wu2026needs}。本文采用记录规模、互动信号和时间新近性作为多源指向度之外的补充维度，但把分数限定为当前数据批次内的应用筛选排序。

自然语言处理深度学习模型综述为字符 TF-IDF、冻结多语言表示与小型 MLP 的模型族比较提供了技术背景\cite{he2025nlp}。本文保留字符模型作为可解释基线，并把 E5 表示和 MLP 放在同一组外层划分中比较，避免把“使用神经网络”直接等同于性能提升。

关于 LLM 标注稳健性的最新研究指出，多目标和高维标签任务可能出现系统性盲点\cite{senthilkumar2026brittle}。这支持本文把 509 条逐条审阅标签上的分类结果限定为复现性检查，并把其余 2,052 条授权转登记字段单独记录为标签来源边界。医疗、电商和风险控制类应用研究则说明，应用方案需要结合领域流程与人工审批，C2 策划案因此保留证据检索、权限隔离和人工确认节点\cite{shi2026tcm,he2024ecommerce}。

\subsection{来源与文本质量审计}

来源结构是多源热点解释的前提。清洗后记录来自 Hacker News 1,523 条、量子位 581 条、GitHub 387 条、V2EX 48 条和 IT 之家 22 条。Hacker News 占据最大比例，因此 C1 的 66 条成员不能独立证明跨平台需求。GitHub、V2EX 和量子位共同出现在 C2 中，C2 的多源属性来自网站覆盖字段，而不是单纯的记录数量。

正文质量标记进一步限定了机会判断的证据强度。清洗记录中 670 条被标为内容不完整，495 条为短文本，907 条没有可用互动指标。一个记录可能同时命中多个标记，缺失互动值没有被替换为零，短文本也没有因为语义相似而自动合并。机会判别要求证据摘录回到清洗正文，因而标题、简介或仅有订阅摘要的记录只能在已有文字范围内判断，不能由模型补写业务背景。

\begin{table}[H]
  \centering
  \small
  \caption{来源规模与主要质量标记}
  \label{tab:data-quality-audit}
  \begin{tabular}{lrrr}
    \toprule
    来源或标记 & 记录数 & 占比（\%） & 解释 \\
    \midrule
    Hacker News & 1523 & 59.47 & 单一来源占比最高 \\
    量子位 & 581 & 22.69 & 中文资讯来源 \\
    GitHub & 387 & 15.11 & 工程项目与 Issue \\
    V2EX & 48 & 1.87 & 中文用户讨论 \\
    IT 之家 & 22 & 0.86 & 订阅入口样本 \\
    内容不完整 & 670 & 26.16 & 可能缺少上下文 \\
    短文本 & 495 & 19.33 & 证据范围较窄 \\
    缺失互动指标 & 907 & 35.42 & 不参与互动维度 \\
    \bottomrule
  \end{tabular}
\end{table}

上述统计的分母是 2,561 条清洗记录；质量标记允许重叠，来源计数互斥。该审计使主题价值分、机会置信度和分类器结果都能回到实际数据字段，避免把平台覆盖差异误写成行业需求比例。
